% Chapter 9: MCP and Tool Integration
\chapter{MCP and Tool Integration}
\label{ch:mcp}

The Model Context Protocol (MCP) extends Claude's capabilities beyond
the built-in tools. It is an open standard that connects Claude to
external services, databases, APIs, and custom tools through a
unified interface.

\section{Model Context Protocol}
\label{sec:mcp-overview}

\marginnote[Official]{Open standard with 10,000+ public servers. Adopted by ChatGPT, Cursor, Gemini, VS Code, Copilot. 100M+ monthly downloads.\sourceurl{https://code.claude.com/docs/en/mcp}}

\official{MCP provides a standardized way to connect Claude to external
tools and data sources.} Key concepts:

\begin{description}
  \item[Servers] Provide tools, resources, and prompts via a standard protocol.
    Transport: HTTP, SSE, or stdio.
  \item[Tools] Functions Claude can call (e.g., query a database, create a ticket).
  \item[Resources] Data Claude can reference (e.g., project docs, API schemas).
    Use with \code{@server:resource/path} syntax.
  \item[Prompts] Pre-built prompt templates provided by the server.
\end{description}

\subsection{Installation and Scopes}

\begin{minted}{bash}
# Add an MCP server
claude mcp add github --transport stdio -- gh-mcp

# Scopes
# Local: current machine only (default)
# Project: .mcp.json committed to git (team-wide)
# User: ~/.claude/mcp.json (personal global)
\end{minted}

\marginnote[Official]{OAuth 2.0 support for authenticated servers. Managed MCP for IT policy enforcement.\sourceurl{https://code.claude.com/docs/en/mcp}}

\subsection{Dynamic Tool Loading}

\official{When tool definitions exceed 10\% of context, Claude uses
dynamic tool loading.} Tools are loaded on-demand via search rather than
all being present in context. This keeps context efficient when using
servers with many tools.

\subsection{Security Considerations}

\begin{warnbox}[MCP Security]
\begin{itemize}
  \item Verify server provenance before installation.
  \item Prefer servers from the official registry.
  \item Review permissions each server requires.
  \item Use Managed MCP for enterprise policy enforcement.
  \item Audit server access logs regularly.
\end{itemize}
\end{warnbox}

% -------------------------------------------------------------------
\section{Practical MCP Patterns}
\label{sec:mcp-patterns}

\subsection{Project-Level MCP Configuration}

\official{Share MCP configs via \code{.mcp.json} committed to git.}
This ensures every team member has the same tool integrations.

\begin{minted}{json}
{
  "mcpServers": {
    "github": {
      "command": "gh-mcp",
      "transport": "stdio"
    },
    "database": {
      "command": "mcp-postgres",
      "args": ["--connection-string", "$DATABASE_URL"],
      "transport": "stdio"
    }
  }
}
\end{minted}

\subsection{Cross-Project Knowledge Bases}

\practitioner{For organizations with multiple repos that share a knowledge base,
use a central MCP server for cross-project queries.}

\begin{enumerate}
  \item Set up a knowledge base MCP server in a shared infrastructure repo.
  \item Each consuming project references it via \code{.mcp.json}.
  \item Claude in any project can query the shared knowledge base.
\end{enumerate}

This enables patterns like:
\begin{itemize}
  \item Querying documentation across all team repos from any project.
  \item Searching for code patterns across the organization.
  \item Accessing shared research and decision records.
\end{itemize}

\subsection{MCP Resources for Context}

\official{Reference MCP resources with \code{@server:resource/path}.}
Resources are auto-fetched when mentioned, providing Claude with
up-to-date external context without manual copy-paste.

\begin{minted}{text}
Prompt: "Using @database:schema/users and @docs:api/authentication,
design a password reset endpoint."
\end{minted}

\begin{keyinsight}
MCP's value is not in any single server---it is in the ecosystem.
A project with GitHub + database + documentation MCP servers enables
Claude to work with the full context of your system, not just the files
on disk.
\end{keyinsight}
